\documentclass[12pt,a4paper]{report}

% ── Encoding & Fonts ──
\usepackage[utf8]{inputenc}
\usepackage[T1]{fontenc}
\usepackage{lmodern}

% ── Page Layout ──
\usepackage[margin=1in,headheight=15pt]{geometry}
\usepackage{setspace}
\onehalfspacing

% ── Language ──
\usepackage[english]{babel}

% ── Graphics & Colors ──
\usepackage{graphicx}
\usepackage[dvipsnames,svgnames,table]{xcolor}

% ── Math ──
\usepackage{amsmath,amssymb,amsfonts}

% ── Tables ──
\usepackage{booktabs}
\usepackage{tabularx}
\usepackage{longtable}
\usepackage{multirow}

% ── Lists ──
\usepackage{enumitem}

% ── Code Listings ──
\usepackage{listings}
\lstset{
  basicstyle=\ttfamily\small,
  keywordstyle=\color{blue}\bfseries,
  commentstyle=\color{gray},
  stringstyle=\color{red},
  breaklines=true,
  frame=single,
  backgroundcolor=\color{gray!5},
  numbers=left,
  numberstyle=\tiny\color{gray},
  numbersep=8pt,
  tabsize=4,
  showstringspaces=false,
}

% ── Hyperlinks ──
\usepackage[colorlinks=true,linkcolor=NavyBlue,citecolor=Green,urlcolor=RoyalBlue]{hyperref}

% ── Headers & Footers ──
\usepackage{fancyhdr}
\pagestyle{fancy}
\fancyhf{}
\fancyhead[L]{\leftmark}
\fancyhead[R]{\thepage}
\renewcommand{\headrulewidth}{0.4pt}

% ── Chapter Style ──
\usepackage{titlesec}
\titleformat{\chapter}[display]
  {\normalfont\huge\bfseries}
  {\chaptertitlename\ \thechapter}{20pt}{\Huge}
\titlespacing*{\chapter}{0pt}{-20pt}{30pt}

% ── Caption Styling ──
\usepackage{caption}
\captionsetup{font=small,labelfont=bf}

% ── Misc ──
\usepackage{parskip}
\usepackage{float}

% ════════════════════════════════════════════════
%  DOCUMENT METADATA
% ════════════════════════════════════════════════
\title{\textbf{Convergent CFD Simulation Guide (v2)}\\[6pt]
       \Large Hydrogen--Methane T-Junction Mixing\\
       Three-Phase Approach in ANSYS Fluent}
\author{Yash}
\date{\today}

\begin{document}

\maketitle
\thispagestyle{empty}
\tableofcontents
\newpage


% ════════════════════════════════════════════════════════════════
%  CHAPTER 1 — Introduction
% ════════════════════════════════════════════════════════════════
\chapter{Introduction}
\label{ch:intro}

\section{Purpose of This Guide}

This document replaces the original \emph{Parameterised CFD Simulation
Guide}.  The original setup launched a transient Spalart--Allmaras DDES
directly from a cold (hybrid-initialised) field.  In practice this
approach frequently \textbf{diverges} in ANSYS Fluent because several
aggressive settings are activated simultaneously:

\begin{enumerate}
  \item \textbf{No steady-state precursor.}  The PISO algorithm
        requires a reasonably established velocity and pressure field
        to function stably; launching it from hybrid initialisation
        introduces large, non-physical transients.
  \item \textbf{Bounded Central Differencing (BCD) from step~1.}
        BCD is low-dissipation --- essential for DES --- but it
        amplifies perturbations on an uninitialised field.
  \item \textbf{Spalart--Allmaras in an internal flow.}  The
        one-equation SA model was designed for external aerodynamics
        and is documented as ``difficult to define length scale in
        complex geometries'' and ``unsuitable for more general
        internal flows'' (Versteeg \& Malalasekera, 2007).
  \item \textbf{High-pressure ideal gas + buoyancy from $t=0$.}  At
        6.9\,MPa, density is hypersensitive to pressure and
        temperature fluctuations.  The full $\rho\,\mathbf{g}$ body
        force creates steep hydrostatic gradients that an
        uninitialised solver struggles to balance.
\end{enumerate}

\subsection{The Fix: Three-Phase Strategy}

The solution is to build the flow field in stages, adding complexity
only when the previous stage has converged.

\begin{table}[H]
\centering
\caption{Three-phase simulation strategy.}
\label{tab:phases}
\begin{tabularx}{\textwidth}{@{}clXl@{}}
\toprule
\textbf{Phase} & \textbf{Mode} & \textbf{Purpose} & \textbf{Duration} \\
\midrule
1 & Steady RANS   & Establish converged pressure, velocity, species
                    and temperature fields & 2\,000--5\,000 iters \\
2 & Transient RANS & Develop the unsteady flow structure (jet
                     shedding, buoyant plume) with stable numerics
                   & $\sim$5\,000 steps (1\,s) \\
3 & Transient DDES & Production run with scale-resolving turbulence;
                     collect time-averaged statistics
                   & 25\,000 steps (5\,s) \\
\bottomrule
\end{tabularx}
\end{table}

\noindent The final Phase~3 results are physically identical to what
the Eames et al.\ (2022) paper produces.  Phases~1 and~2 are
\emph{numerical preparation} --- they do not alter the physics.


\section{Problem Definition}

Pure hydrogen is injected via a $90^{\circ}$ T-junction into a
horizontal pipeline carrying methane (modelling natural gas) at
6.9\,MPa.  The density ratio $\rho_{\mathrm{CH}_4}/\rho_{\mathrm{H}_2}
\approx 8$ generates buoyancy-driven stratification and elevated wall
H$_2$ concentrations relevant to embrittlement risk.

\begin{table}[H]
\centering
\caption{Reference conditions (Eames et al., 2022).}
\label{tab:refcond}
\begin{tabular}{@{}llr@{}}
\toprule
\textbf{Parameter} & \textbf{Symbol} & \textbf{Value} \\
\midrule
Main-pipe diameter   & $d_1$           & 460\,mm \\
Branch diameter      & $d_2$           & 115\,mm ($\beta = 0.25$) \\
Operating pressure   & $p_1$           & 6.9\,MPa \\
Main-flow speed      & $u_1$           & 10\,m/s \\
Branch speed         & $u_2$           & 32\,m/s ($\mathcal{D} = 0.2$) \\
Main temperature     & $T_1$           & 284\,K \\
Branch temperature   & $T_2$           & 293\,K \\
Injection angle      & $\theta$        & $90^{\circ}$ (top-side) \\
Main-pipe length     & ---             & $20\,d_1 = 9200$\,mm \\
Junction location    & ---             & $10\,d_1 = 4600$\,mm from inlet \\
\bottomrule
\end{tabular}
\end{table}


\section{Turbulence Model Choice: $k$--$\omega$ SST}
\label{sec:whysst}

The original guide used \textbf{Spalart--Allmaras (SA)}.  This guide
switches to \textbf{$k$--$\omega$ SST} (Menter, 1994) for two
reasons:

\begin{enumerate}
  \item \textbf{Better suited to internal flows.}  SST blends near-wall
        $k$--$\omega$ behaviour (accurate in adverse pressure gradients,
        separation and impingement) with free-stream $k$--$\varepsilon$
        behaviour (avoids the inlet-$\omega$ sensitivity problem).
        Versteeg \& Malalasekera rate SST as ``most general'' among
        one- and two-equation models.
  \item \textbf{SST-based DDES is available in Fluent.}  When you
        enable DES on an SST model, Fluent uses the SST-DDES
        formulation, which provides better shielding of the boundary
        layer (the ``delayed'' switch) and better resolution of
        separated shear layers than SA-DDES in internal geometries.
\end{enumerate}

\noindent The paper (Eames 2022) used SA-DES in OpenFOAM.  Using
SST-DDES is an \emph{improvement}, not a departure --- the flow
physics (buoyant mixing, shedding, stratification) are captured
equally well, with better convergence behaviour and wall-bounded
accuracy.


% ════════════════════════════════════════════════════════════════
%  CHAPTER 2 — Geometry and Meshing
% ════════════════════════════════════════════════════════════════
\chapter{Geometry and Meshing}
\label{ch:mesh}

The geometry and meshing procedure is \textbf{identical} to the
original guide (Chapters~2--3).  Only a brief summary is provided
here; refer to the original document for step-by-step DesignModeler
and ANSYS Meshing instructions.

\section{Geometry Summary}

\begin{itemize}
  \item Main pipe: 460\,mm diameter, 9200\,mm length along $Z$.
  \item Branch: 115\,mm diameter at $Z = 4600$\,mm, injection
        along $+Y$ (top-side).
  \item Five Named Selections: \texttt{main\_inlet},
        \texttt{branch\_inlet}, \texttt{outlet}, \texttt{wall},
        \texttt{BOI\_Body}.
\end{itemize}

\section{Meshing Summary}

\begin{itemize}
  \item Poly-hexcore method on \texttt{Fluid\_Domain}.
  \item BOI element size: 5\,mm (within 690\,mm-radius sphere at
        junction).
  \item 20 inflation layers on \texttt{wall}, first layer
        height 0.001\,mm ($y^+ \approx 1$), growth rate 1.2.
  \item Target cell count: $\sim$3\,million.
\end{itemize}

\section{Pre-Flight Mesh Check}

Before proceeding to the solver, verify:

\begin{enumerate}
  \item \textbf{Orthogonal Quality} $> 0.1$ (check in Meshing
        $\rightarrow$ Statistics $\rightarrow$ Mesh Metric).
  \item \textbf{Aspect Ratio} $< 100$ (except in inflation layers).
  \item \textbf{Skewness} $< 0.95$ for all cells.
  \item \textbf{Named Selections} all present and correctly scoped.
\end{enumerate}

\noindent\fbox{\parbox{\dimexpr\textwidth-2\fboxsep-2\fboxrule}{%
\textbf{If the mesh has quality problems:}  Poor orthogonality or
high skewness are the \#1 hidden cause of divergence.  Fix the mesh
before touching solver settings.  A converged solution on a poor mesh
is worse than a divergent run on a good mesh --- at least the
divergence tells you something is wrong.}}


% ════════════════════════════════════════════════════════════════
%  CHAPTER 3 — Fluent Launch and Common Settings
% ════════════════════════════════════════════════════════════════
\chapter{Fluent Launch and Settings Common to All Phases}
\label{ch:common}

\section{Launching Fluent}

Double-click the \textbf{Setup} cell in Workbench.  When prompted:
\begin{itemize}
  \item \textbf{Double Precision:} Yes.
  \item \textbf{Solver Processes:} 8 (for an i7-14700 with 8 P-cores).
\end{itemize}

\section{Materials (All Phases)}

The material setup is \textbf{identical} across all three phases.

\subsection{Phase 1: Import Pure Gases}

\begin{enumerate}
  \item \textbf{Materials} $\rightarrow$ \textbf{Fluid} $\rightarrow$
        \textbf{Create/Edit}.
  \item Click \textbf{Fluent Database}.
  \item Search for \texttt{methane (ch4)} $\rightarrow$ \textbf{Copy},
        then \textbf{Close}.
  \item Repeat for \texttt{hydrogen (h2)}.
\end{enumerate}

\subsection{Phase 2: Create Custom Mixture}

\begin{enumerate}
  \item In the Materials panel, click \textbf{Create/Edit} under
        \textbf{Mixture}.
  \item Name: \texttt{ch4-h2-mixture}.
  \item \textbf{Species:} Add \texttt{h2} and \texttt{ch4}; remove
        \texttt{o2} and \texttt{n2}.
  \item \textbf{Order:} \texttt{h2} first, \texttt{ch4} second
        (CH$_4$ is the balance species).
\end{enumerate}

\subsection{Phase 3: Mixture Properties}

\begin{table}[H]
\centering
\caption{Mixture-level property models.}
\begin{tabular}{@{}ll@{}}
\toprule
\textbf{Property} & \textbf{Model} \\
\midrule
Density              & \texttt{ideal-gas} \\
$c_p$                & \texttt{mixing-law} \\
Thermal conductivity & \texttt{ideal-gas-mixing-law} \\
Viscosity            & \texttt{ideal-gas-mixing-law} \\
Mass diffusivity     & \texttt{kinetic-theory} \\
\bottomrule
\end{tabular}
\end{table}

\subsection{NIST Verification}

\begin{table}[H]
\centering
\caption{Verify individual species properties against NIST (at
$T \approx 284$\,K, $p = 6.9$\,MPa).}
\begin{tabular}{@{}lrrl@{}}
\toprule
\textbf{Property} & \textbf{CH$_4$} & \textbf{H$_2$} & \textbf{Source} \\
\midrule
Molecular weight (kg\,kmol$^{-1}$) & 16.043 & 2.016 & NIST \\
$c_p$ (J\,kg$^{-1}$\,K$^{-1}$)    & 2580   & 14\,500 & NIST \\
Thermal conductivity (W\,m$^{-1}$\,K$^{-1}$) & 0.038 & 0.185 & NIST \\
Dynamic viscosity (Pa\,s) & $1.25 \times 10^{-5}$ & $9.0 \times 10^{-6}$ & NIST \\
\bottomrule
\end{tabular}
\end{table}


\section{Operating Conditions (All Phases)}

\begin{enumerate}
  \item \textbf{Operating Pressure:} \texttt{6900000}\,Pa.
  \item \textbf{Reference Pressure Location:} Default $(0, 0, 0)$.
  \item \textbf{Gravity:} Enable.  $g_x = 0$, $g_y = -9.81$,
        $g_z = 0$.
  \item \textbf{Specified Operating Density:} Enable, value
        \textbf{0}.
\end{enumerate}

\noindent\fbox{\parbox{\dimexpr\textwidth-2\fboxsep-2\fboxrule}{%
\textbf{Why operating density = 0?}  With an ideal-gas compressible
mixture and gravity enabled, the buoyancy force must be
$\rho\,\mathbf{g}$ (using the true local density), not
$(\rho-\rho_0)\mathbf{g}$.  A non-zero $\rho_0$ introduces an
artificial hydrostatic baseline that can destabilise the
calculation.}}


\section{Species Model (All Phases)}

\begin{enumerate}
  \item \textbf{Models} $\rightarrow$ \textbf{Species} $\rightarrow$
        \textbf{Species Transport}.
  \item \textbf{Reactions:} Disabled.
  \item \textbf{Diffusion Energy Source:} Enabled.
  \item \textbf{Full Multicomponent Diffusion:} Disabled.
\end{enumerate}


\section{Energy Equation (All Phases)}

\textbf{Enable} the Energy Equation.

\section{Turbulent Schmidt Number}

Set $Sc_t = 0.5$ via TUI:

\begin{lstlisting}[language={},numbers=none]
/define/models/viscous/turbulent-schmidt-number 0.5
\end{lstlisting}


\section{Boundary Conditions (All Phases)}

The boundary values are \textbf{identical} across all phases.  Only
the solver mode and discretisation change.

\subsection{\texttt{main\_inlet} --- Velocity Inlet}

\begin{table}[H]
\centering
\caption{\texttt{main\_inlet} settings.}
\begin{tabular}{@{}llr@{}}
\toprule
\textbf{Tab} & \textbf{Setting} & \textbf{Value} \\
\midrule
Momentum & Velocity Magnitude     & 10\,m/s \\
Momentum & Turbulence Spec.       & Intensity \& Hydraulic Diameter \\
Momentum & Turbulent Intensity    & 5\,\% \\
Momentum & Hydraulic Diameter     & 0.46\,m \\
Thermal  & Temperature            & 284\,K \\
Species  & Mass Fraction (h2)     & 0 \\
\bottomrule
\end{tabular}
\end{table}

\subsection{\texttt{branch\_inlet} --- Velocity Inlet}

\begin{table}[H]
\centering
\caption{\texttt{branch\_inlet} settings.}
\begin{tabular}{@{}llr@{}}
\toprule
\textbf{Tab} & \textbf{Setting} & \textbf{Value} \\
\midrule
Momentum & Velocity Magnitude     & 32\,m/s \\
Momentum & Turbulence Spec.       & Intensity \& Hydraulic Diameter \\
Momentum & Turbulent Intensity    & 5\,\% \\
Momentum & Hydraulic Diameter     & 0.115\,m \\
Thermal  & Temperature            & 293\,K \\
Species  & Mass Fraction (h2)     & 1 \\
\bottomrule
\end{tabular}
\end{table}

\noindent\textbf{Branch velocity derivation:}
$\mathcal{D} = \beta^2 (u_2/u_1)$ gives
$0.2 = (0.25)^2 \times u_2/10$, hence $u_2 = 32$\,m/s.

\subsection{\texttt{outlet} --- Pressure Outlet}

\begin{table}[H]
\centering
\caption{\texttt{outlet} settings.}
\begin{tabular}{@{}llr@{}}
\toprule
\textbf{Tab} & \textbf{Setting} & \textbf{Value} \\
\midrule
Momentum & Gauge Pressure            & 0\,Pa \\
Momentum & Backflow Turbulence Spec. & Intensity \& Hydraulic Diameter \\
Momentum & Backflow Turb.\ Intensity & 5\,\% \\
Momentum & Backflow Hydraulic Dia.   & 0.46\,m \\
Thermal  & Backflow Temperature      & 284\,K \\
Species  & Backflow h2 Mass Frac.    & 0 \\
\bottomrule
\end{tabular}
\end{table}

\subsection{\texttt{wall} --- Wall}

\begin{tabular}{@{}ll@{}}
\toprule
\textbf{Tab} & \textbf{Setting} \\
\midrule
Momentum & Stationary Wall, No Slip \\
Thermal  & Heat Flux = 0\,W/m$^2$ (adiabatic) \\
\bottomrule
\end{tabular}


% ════════════════════════════════════════════════════════════════
%  CHAPTER 4 — Phase 1: Steady-State RANS
% ════════════════════════════════════════════════════════════════
\chapter{Phase 1: Steady-State RANS}
\label{ch:phase1}

The goal of this phase is to produce a \textbf{converged, smooth
velocity--pressure--species field} that the transient solver can use
as a starting point.  No time-dependent physics are captured here ---
this is purely numerical preparation.


\section{Step 1: General Settings}

\begin{enumerate}
  \item \textbf{Type:} Pressure-Based.
  \item \textbf{Velocity Formulation:} Absolute.
  \item \textbf{Time:} \textbf{Steady}.
\end{enumerate}

\noindent Gravity and operating conditions were already set in
Chapter~\ref{ch:common}.


\section{Step 2: Turbulence Model}

\begin{enumerate}
  \item \textbf{Models} $\rightarrow$ \textbf{Viscous}.
  \item Select \textbf{$k$--omega (2\,eqn)}.
  \item \textbf{$k$--omega Model:} \textbf{SST}.
  \item Leave all \textbf{Model Constants} at defaults.
  \item Do \textbf{NOT} enable DES yet.
\end{enumerate}

\noindent\fbox{\parbox{\dimexpr\textwidth-2\fboxsep-2\fboxrule}{%
\textbf{Why SST and not SA?}  SST is a two-equation model that
handles internal flows with separation, impingement and curved shear
layers far better than the one-equation SA model.  SST is also the
natural base for the DDES formulation we will enable in Phase~3.}}


\section{Step 3: Solution Methods}

\begin{table}[H]
\centering
\caption{Phase~1 solution methods.}
\begin{tabular}{@{}ll@{}}
\toprule
\textbf{Setting} & \textbf{Value} \\
\midrule
Pressure--Velocity Coupling  & \textbf{Coupled} \\
\midrule
\multicolumn{2}{@{}l}{\textbf{Spatial Discretisation}} \\
\midrule
Gradient                     & Least Squares Cell Based \\
Pressure                     & Second Order \\
Momentum                     & \textbf{Second Order Upwind} \\
Turbulent Kinetic Energy     & Second Order Upwind \\
Specific Dissipation Rate    & Second Order Upwind \\
Species                      & Second Order Upwind \\
Energy                       & Second Order Upwind \\
\bottomrule
\end{tabular}
\end{table}

\noindent\fbox{\parbox{\dimexpr\textwidth-2\fboxsep-2\fboxrule}{%
\textbf{Why Coupled, not PISO?}  The Coupled solver solves the
pressure and momentum equations simultaneously.  For steady-state
problems it converges \textbf{5--10$\times$ faster} than segregated
SIMPLE/SIMPLEC.  It also handles compressible buoyant flows more
robustly because the pressure--density coupling is tighter.\\[6pt]
\textbf{Why Second Order Upwind, not BCD?}  Upwind schemes add
numerical dissipation which is \emph{bad} for resolving turbulent
eddies (Phase~3) but \emph{good} for convergence.  We only need a
smooth field here --- accuracy of vortex structures does not matter
in steady RANS.}}


\section{Step 4: Under-Relaxation Factors}

The Coupled solver uses \textbf{Courant Number} instead of separate
pressure/momentum URFs.

\begin{table}[H]
\centering
\caption{Phase~1 solution controls.}
\begin{tabular}{@{}lr@{}}
\toprule
\textbf{Control} & \textbf{Value} \\
\midrule
Courant Number (Coupled)  & \textbf{50} \\
\midrule
\multicolumn{2}{@{}l}{\textbf{Explicit Relaxation Factors}} \\
\midrule
Momentum                  & 0.5 \\
Pressure                  & 0.5 \\
\midrule
\multicolumn{2}{@{}l}{\textbf{Under-Relaxation Factors}} \\
\midrule
Density                   & 0.8 \\
Body Forces               & 0.8 \\
Turbulent Kinetic Energy  & 0.7 \\
Specific Dissipation Rate & 0.7 \\
Turbulent Viscosity       & 0.8 \\
Species                   & 0.8 \\
Energy                    & 0.9 \\
\bottomrule
\end{tabular}
\end{table}

\noindent These are deliberately conservative.  If the solution
converges smoothly, you can increase the Courant Number to 100--200
after the first few hundred iterations to speed up convergence.

\noindent\fbox{\parbox{\dimexpr\textwidth-2\fboxsep-2\fboxrule}{%
\textbf{If it still oscillates:}  Drop the Courant Number to
\textbf{20} and reduce momentum/pressure explicit relaxation to
\textbf{0.25}.  This is very slow but will always converge.  Once
residuals start dropping, gradually increase.}}


\section{Step 5: Initialisation}

\begin{enumerate}
  \item Select \textbf{Standard Initialization} (not Hybrid).
  \item \textbf{Compute From:} \texttt{main\_inlet}.
  \item Click \textbf{Initialize}.
\end{enumerate}

\noindent Standard initialisation from the dominant inlet gives a
uniform-flow starting field that is much gentler on the Coupled solver
than Hybrid initialisation.

\subsection{Patch the Branch}

\begin{enumerate}
  \item Click \textbf{Patch\ldots}
  \item Create a cell register (cylinder or sphere) enclosing the
        branch pipe volume.
  \item Patch \textbf{H$_2$ mass fraction = 1.0}.
  \item Patch \textbf{Y-velocity = $-32$\,m/s} (branch direction).
\end{enumerate}


\section{Step 6: Residual Monitors}

\begin{enumerate}
  \item Set convergence criteria to \textbf{1e-4} for all equations.
  \item Set energy convergence to \textbf{1e-6}.
  \item Enable \textbf{Plot}.
\end{enumerate}


\section{Step 7: Run Phase 1}

\begin{enumerate}
  \item \textbf{Number of Iterations:} 3000.
  \item Click \textbf{Calculate}.
\end{enumerate}

\subsection{What to Watch}

\begin{table}[H]
\centering
\caption{Phase~1 convergence targets.}
\begin{tabular}{@{}lll@{}}
\toprule
\textbf{Indicator} & \textbf{Healthy} & \textbf{Problem} \\
\midrule
Continuity residual & $< 10^{-3}$ & Oscillating or flat \\
Momentum residuals  & $< 10^{-4}$ & Rising \\
Species residual    & $< 10^{-4}$ & Flat above $10^{-2}$ \\
Energy residual     & $< 10^{-6}$ & Above $10^{-3}$ \\
Mass flow balance   & $< 0.5\,\%$ imbalance & $> 2\,\%$ \\
\bottomrule
\end{tabular}
\end{table}

\noindent If residuals flatten but do not reach the targets, this may
indicate low-frequency unsteadiness (e.g.\ vortex shedding at the
junction).  That is \textbf{expected and acceptable} --- the purpose
of Phase~1 is to get a ``reasonable'' field, not a fully converged
steady state.  Proceed to Phase~2 once residuals have dropped by at
least \textbf{3 orders of magnitude} from initialisation.


\section{Step 8: Save}

\begin{enumerate}
  \item \textbf{File} $\rightarrow$ \textbf{Write} $\rightarrow$
        \textbf{Case \& Data\ldots}
  \item Name: \texttt{phase1\_steady\_rans.cas.h5}
\end{enumerate}

\noindent This is your safety checkpoint.  If Phase~2 diverges, you
can reload this file and adjust settings without re-running Phase~1.


% ════════════════════════════════════════════════════════════════
%  CHAPTER 5 — Phase 2: Transient RANS
% ════════════════════════════════════════════════════════════════
\chapter{Phase 2: Transient RANS}
\label{ch:phase2}

Phase~2 introduces time-dependence while keeping the robust SST RANS
turbulence model.  The goal is to establish the \textbf{unsteady flow
structure} --- the jet shedding, buoyant plume, and recirculation
zones --- using numerics that can handle the transient without
diverging.


\section{Step 1: Switch to Transient}

\begin{enumerate}
  \item \textbf{Setup} $\rightarrow$ \textbf{General} $\rightarrow$
        \textbf{Time:} change from Steady to \textbf{Transient}.
\end{enumerate}

\noindent\fbox{\parbox{\dimexpr\textwidth-2\fboxsep-2\fboxrule}{%
\textbf{Do NOT change the turbulence model yet.}  Keep
$k$--$\omega$ SST.  The DES switch comes in Phase~3, after the
transient flow structure is established.}}


\section{Step 2: Change Pressure--Velocity Coupling}

\begin{enumerate}
  \item \textbf{Solution} $\rightarrow$ \textbf{Methods}.
  \item Change \textbf{Pressure--Velocity Coupling} from Coupled to
        \textbf{PISO}.
  \item Skewness Correction: \textbf{1} (default).
  \item Neighbour Correction: \textbf{1} (default).
\end{enumerate}


\section{Step 3: Update Spatial Discretisation}

\begin{table}[H]
\centering
\caption{Phase~2 spatial discretisation (changes from Phase~1
         highlighted).}
\begin{tabular}{@{}lll@{}}
\toprule
\textbf{Field} & \textbf{Phase~1} & \textbf{Phase~2} \\
\midrule
Gradient   & Least Sq.\ Cell Based & \emph{same} \\
Pressure   & Second Order          & \textbf{PRESTO!} \\
Momentum   & Second Order Upwind   & \textbf{Bounded Central Diff.} \\
$k$        & Second Order Upwind   & \emph{same} \\
$\omega$   & Second Order Upwind   & \emph{same} \\
Species    & Second Order Upwind   & \emph{same} \\
Energy     & Second Order Upwind   & \emph{same} \\
\bottomrule
\end{tabular}
\end{table}

\noindent\fbox{\parbox{\dimexpr\textwidth-2\fboxsep-2\fboxrule}{%
\textbf{Why switch momentum to BCD now?}  BCD is essential for
resolving turbulent structures in Phase~3 (DDES).  Introducing it
in Phase~2 on an already-converged field avoids the oscillations that
BCD causes on an uninitialised field.  If BCD causes residual spikes
here, temporarily revert to \textbf{Second Order Upwind} for 500
steps, then switch back.}}


\section{Step 4: Transient Formulation}

\begin{enumerate}
  \item \textbf{Transient Formulation:} \textbf{Bounded Second Order
        Implicit}.
\end{enumerate}


\section{Step 5: Under-Relaxation Factors}

Switch from the Coupled-solver controls to PISO-style URFs:

\begin{table}[H]
\centering
\caption{Phase~2 under-relaxation factors.}
\begin{tabular}{@{}lr@{}}
\toprule
\textbf{Variable} & \textbf{Factor} \\
\midrule
Pressure                     & 0.3 \\
Momentum                     & 0.7 \\
Turbulent Kinetic Energy     & 0.8 \\
Specific Dissipation Rate    & 0.8 \\
Turbulent Viscosity          & 0.9 \\
Species                      & 0.9 \\
Energy                       & 1.0 \\
Density                      & 1.0 \\
Body Forces                  & 1.0 \\
\bottomrule
\end{tabular}
\end{table}


\section{Step 6: Time Stepping --- Conservative Start}

\begin{table}[H]
\centering
\caption{Phase~2 time stepping.}
\begin{tabular}{@{}llr@{}}
\toprule
\textbf{Parameter} & \textbf{First 1000 steps} & \textbf{After 1000 steps} \\
\midrule
Time Step Size (s) & \textbf{5e-4}           & \textbf{2e-4} \\
CFL (approx.)      & $\sim 0.6$              & $\sim 1.3$ \\
Max Iters / Step   & 30                      & 20 \\
\bottomrule
\end{tabular}
\end{table}

\begin{enumerate}
  \item Set \textbf{Time Step Size} to \texttt{5e-4}.
  \item Set \textbf{Number of Time Steps} to \texttt{1000}.
  \item Set \textbf{Max Iterations / Time Step} to \texttt{30}.
  \item Click \textbf{Calculate}.
  \item After 1000 steps complete, change Time Step Size to
        \texttt{2e-4}.
  \item Set Number of Time Steps to \texttt{4000}.
  \item Set Max Iterations / Time Step to \texttt{20}.
  \item Click \textbf{Calculate}.
\end{enumerate}

\noindent Total Phase~2 time: $1000 \times 5\!\times\!10^{-4}
+ 4000 \times 2\!\times\!10^{-4} = 0.5 + 0.8 = 1.3$\,s of physical
time.

\subsection{What to Watch}

\begin{itemize}
  \item \textbf{Residuals:} Should drop 1--2 orders of magnitude
        \emph{within each time step}.  If they plateau at
        $\sim 10^{-2}$ and oscillate, this is acceptable for PISO.
  \item \textbf{Max Courant Number:} Printed in the console.
        Target: $< 5$.  If consistently above 10, reduce $\Delta t$.
  \item \textbf{Outlet H$_2$:} Set up a surface monitor (mass-weighted
        average of h2 on outlet).  It should start developing toward
        $\sim 0.019$ mass fraction.
  \item \textbf{Negative temperatures:} If these appear, immediately
        reduce $\Delta t$ to \texttt{1e-4} and run 500 steps.
\end{itemize}


\section{Step 7: Save}

\begin{enumerate}
  \item \textbf{File} $\rightarrow$ \textbf{Write} $\rightarrow$
        \textbf{Case \& Data\ldots}
  \item Name: \texttt{phase2\_transient\_rans.cas.h5}
\end{enumerate}


% ════════════════════════════════════════════════════════════════
%  CHAPTER 6 — Phase 3: Transient DDES (Production)
% ════════════════════════════════════════════════════════════════
\chapter{Phase 3: Transient DDES (Production Run)}
\label{ch:phase3}

This is the final phase that produces publication-quality results.
The flow field from Phase~2 provides a stable starting point for the
scale-resolving DDES turbulence model.


\section{Step 1: Enable DDES}

\begin{enumerate}
  \item \textbf{Models} $\rightarrow$ \textbf{Viscous}.
  \item The model should still show \textbf{$k$--omega (2\,eqn), SST}.
  \item In the left panel, select \textbf{Detached Eddy Simulation
        (DES)}.
  \item \textbf{DES Option:} \textbf{Delayed DES (DDES)}.
  \item Leave all model constants at defaults.
  \item Click \textbf{OK}.
\end{enumerate}

\noindent This changes the turbulence model from RANS to SST-DDES.
In attached boundary layers, the solver continues using the SST RANS
equations.  In separated/free-shear regions (the junction wake, the
buoyant plume), it switches to LES-like behaviour, resolving large
eddies directly.

\noindent\fbox{\parbox{\dimexpr\textwidth-2\fboxsep-2\fboxrule}{%
\textbf{Alternative: SA-DDES.}  If you specifically want to replicate
the Eames 2022 paper, you can switch from SST to Spalart--Allmaras
and enable DES$\rightarrow$DDES.  However, SST-DDES is generally
more robust and accurate for internal flows.  Both will produce
qualitatively the same stratification and mixing results.}}


\section{Step 2: Verify All Other Settings}

\textbf{Nothing else changes from Phase~2.}  Confirm:

\begin{table}[H]
\centering
\caption{Phase~3 settings checklist.}
\begin{tabular}{@{}lll@{}}
\toprule
\textbf{Setting} & \textbf{Required Value} & \textbf{Check} \\
\midrule
Time                   & Transient             & $\checkmark$ \\
P--V Coupling          & PISO                  & $\checkmark$ \\
Momentum Scheme        & Bounded Central Diff. & $\checkmark$ \\
Pressure Scheme        & PRESTO!               & $\checkmark$ \\
Transient Formulation  & Bounded 2nd Order Implicit & $\checkmark$ \\
Time Step Size         & 2e-4\,s               & $\checkmark$ \\
Max Iters / Step       & 20                    & $\checkmark$ \\
\bottomrule
\end{tabular}
\end{table}


\section{Step 3: Set Up Monitors}

\subsection{Residual Monitors}

\begin{enumerate}
  \item Equations: continuity, x/y/z-velocity, $k$, $\omega$ (or
        $\tilde{\nu}$ if using SA), h2, energy.
  \item \textbf{Convergence Criterion:} \texttt{1e-4} for all;
        energy \texttt{1e-6}.
  \item \textbf{Plot:} Enabled.
\end{enumerate}

\subsection{Surface Monitors}

\begin{table}[H]
\centering
\caption{Surface monitors.}
\begin{tabular}{@{}lllr@{}}
\toprule
\textbf{Name} & \textbf{Type} & \textbf{Surface} & \textbf{Field} \\
\midrule
\texttt{outlet\_h2}   & Mass-Wtd Average & outlet      & h2 mass frac. \\
\texttt{outlet\_mdot} & Mass Flow Rate   & outlet      & --- \\
\texttt{main\_mdot}   & Mass Flow Rate   & main\_inlet & --- \\
\texttt{branch\_mdot} & Mass Flow Rate   & branch\_inlet & --- \\
\bottomrule
\end{tabular}
\end{table}

\subsection{Point Monitors (Wall H$_2$)}

\begin{table}[H]
\centering
\caption{Wall probe locations ($d_1 = 0.46$\,m, junction at $Z = 4.6$\,m).}
\begin{tabular}{@{}lrrrr@{}}
\toprule
\textbf{Name} & $X$ (m) & $Y$ (m) & $Z$ (m) & $z/d_1$ \\
\midrule
\texttt{top\_z2d}    & 0 &  0.23 & 5.52 & 2 \\
\texttt{top\_z4d}    & 0 &  0.23 & 6.44 & 4 \\
\texttt{top\_z6d}    & 0 &  0.23 & 7.36 & 6 \\
\texttt{top\_z8d}    & 0 &  0.23 & 8.28 & 8 \\
\texttt{bot\_z2d}    & 0 & $-0.23$ & 5.52 & 2 \\
\texttt{bot\_z4d}    & 0 & $-0.23$ & 6.44 & 4 \\
\texttt{bot\_z6d}    & 0 & $-0.23$ & 7.36 & 6 \\
\texttt{bot\_z8d}    & 0 & $-0.23$ & 8.28 & 8 \\
\bottomrule
\end{tabular}
\end{table}

\noindent Monitor both top and bottom walls to capture the
stratification asymmetry that is the central result of the paper.


\section{Step 4: Autosave}

\begin{enumerate}
  \item \textbf{Calculation Activities} $\rightarrow$
        \textbf{Autosave Every:} \texttt{500} Time Steps.
  \item Save both \textbf{Case} and \textbf{Data}.
\end{enumerate}


\section{Step 5: Run Calculation}

\begin{table}[H]
\centering
\caption{Phase~3 run parameters.}
\begin{tabular}{@{}lr@{}}
\toprule
\textbf{Parameter} & \textbf{Value} \\
\midrule
Time Step Size (s)         & \texttt{2e-4} \\
Number of Time Steps       & \texttt{25000} \\
Max Iterations / Time Step & \texttt{20} \\
\midrule
Total physical time        & 5\,s \\
\bottomrule
\end{tabular}
\end{table}

\begin{enumerate}
  \item Click \textbf{Calculate}.
\end{enumerate}


\section{Step 6: Enable Time Statistics (After $t = 4$\,s)}

After \textbf{20\,000} time steps ($t = 4$\,s), the initial transient
has washed out.  Enable time-averaging:

\begin{enumerate}
  \item \textbf{Solution} $\rightarrow$ \textbf{Run Calculation}
        $\rightarrow$ \textbf{Data Sampling for Time Statistics}.
  \item \textbf{Enable:} checked.
  \item \textbf{Sampling Interval:} 1 (every time step).
  \item Continue the remaining 5\,000 steps to $t = 5$\,s.
\end{enumerate}

\noindent This matches the paper's approach: average over $t = 4$--5\,s.


\section{Step 7: What to Watch During the Run}

\begin{table}[H]
\centering
\caption{Runtime diagnostics.}
\begin{tabular}{@{}llll@{}}
\toprule
\textbf{Indicator} & \textbf{Healthy} & \textbf{Warning} & \textbf{Action} \\
\midrule
Residuals per step & Drop 1--2 orders & Flat at $10^{-1}$ & Reduce $\Delta t$ \\
Max Courant        & $< 3$ & $> 5$ consistently & Halve $\Delta t$ \\
Outlet H$_2$       & Stabilise $\sim 0.019$ & Drift $> 0.03$ & Check BCs \\
Mass imbalance     & $< 1\,\%$ & $> 2\,\%$ & Check outlet BC \\
Negative $T$       & Never & Any occurrence & Halve $\Delta t$ \\
\bottomrule
\end{tabular}
\end{table}


\section{Step 8: Save}

\begin{enumerate}
  \item \textbf{File} $\rightarrow$ \textbf{Write} $\rightarrow$
        \textbf{Case \& Data\ldots}
  \item Name: \texttt{phase3\_ddes\_final.cas.h5}
\end{enumerate}


% ════════════════════════════════════════════════════════════════
%  CHAPTER 7 — Post-Processing and Validation
% ════════════════════════════════════════════════════════════════
\chapter{Post-Processing and Validation}
\label{ch:post}

\section{Hydrogen Volume Fraction}

The paper reports $\phi_{\mathrm{H}_2}$ (volume fraction), not
$Y_{\mathrm{H}_2}$ (mass fraction).  Define a Custom Field Function:

\begin{equation}
  \phi_{\mathrm{H}_2}
  = \frac{Y_{\mathrm{H}_2}/M_{\mathrm{H}_2}}
         {Y_{\mathrm{H}_2}/M_{\mathrm{H}_2}
          + (1-Y_{\mathrm{H}_2})/M_{\mathrm{CH}_4}}
\end{equation}

\noindent In Fluent: \textbf{Custom Field Functions} $\rightarrow$
\textbf{Define}:

\begin{lstlisting}[language={},numbers=none]
(Y_H2/2.016) / (Y_H2/2.016 + (1-Y_H2)/16.043)
\end{lstlisting}


\section{Wall $y^+$ Check}

\begin{enumerate}
  \item \textbf{Results} $\rightarrow$ \textbf{Plots} $\rightarrow$
        \textbf{XY Plot}.
  \item Surface: \texttt{wall}.
  \item Field: Wall Fluxes $\rightarrow$ Wall Yplus.
\end{enumerate}

\noindent Target: $y^+ < 5$ everywhere.  Localised $y^+ \sim 2$--3
near stagnation is acceptable.  If $y^+ > 5$ consistently, refine
the inflation layer and re-run.


\section{Expected Results ($\mathcal{D} = 0.2$, $\beta = 0.25$)}

\begin{table}[H]
\centering
\caption{Validation targets (Eames 2022).}
\begin{tabular}{@{}lr@{}}
\toprule
\textbf{Quantity} & \textbf{Expected Value} \\
\midrule
Outlet $V_R$ (volume fraction) & $\approx 0.149$ \\
Outlet $Y_{\mathrm{H}_2}$ (mass fraction) & $\approx 0.019$ \\
Mass balance residual & $< 1\,\%$ \\
Stratification: top wall $V_R$ at $z/d_1 = 8$ & $\approx 1.6 \times V_{R,T}$ \\
Top wall $>$ bottom wall H$_2$ & Yes (buoyancy) \\
\bottomrule
\end{tabular}
\end{table}


\section{Coefficient of Variation (COV)}

On any cross-section downstream of the junction:

\[
  \text{COV} = \frac{\sigma_{\phi}}{\bar{\phi}}
\]

\noindent $\text{COV} = 0$: perfect mixing.
$\text{COV} = 1$: complete segregation.


\section{Mesh Independence (GCI)}

Run three meshes (coarse, medium, fine) and compute:

\[
  p = \frac{\ln|(\phi_3 - \phi_2)/(\phi_2 - \phi_1)|}{\ln r},
  \qquad
  \text{GCI}_{\text{fine}}
  = \frac{1.25\,|(\phi_2 - \phi_1)/\phi_1|}{r^p - 1}
\]

\noindent Target: $\text{GCI} < 2\,\%$.


% ════════════════════════════════════════════════════════════════
%  CHAPTER 8 — Troubleshooting
% ════════════════════════════════════════════════════════════════
\chapter{Troubleshooting}
\label{ch:trouble}

\section{Phase 1 Won't Converge}

\begin{table}[H]
\centering
\caption{Phase~1 troubleshooting.}
\begin{tabularx}{\textwidth}{@{}lXX@{}}
\toprule
\textbf{Symptom} & \textbf{Cause} & \textbf{Fix} \\
\midrule
Residuals oscillate wildly &
  Courant Number too high for the mesh &
  Reduce to \textbf{20}, then increase gradually \\
Continuity residual stuck at $10^{-1}$ &
  Mesh quality issue (high skewness/aspect ratio near junction) &
  Check \textbf{Mesh Quality} in Meshing; fix problem cells \\
Reversed flow at outlet &
  Outlet too close to junction, or backflow BC wrong &
  Ensure outlet is at $Z = 9.2$\,m ($10\,d_1$ from junction) \\
Negative temperatures &
  Operating pressure wrong (not 6.9\,MPa), or energy equation URF
  too high &
  Verify operating pressure = 6\,900\,000; reduce energy URF to 0.7 \\
\bottomrule
\end{tabularx}
\end{table}


\section{Phase 2 Diverges on First Step}

This almost always means the Phase~1 field has a problem:

\begin{enumerate}
  \item Reload \texttt{phase1\_steady\_rans.cas.h5}.
  \item Run 500 more steady iterations.
  \item Try Phase~2 again with $\Delta t = 1\!\times\!10^{-4}$\,s
        (half the recommended value).
\end{enumerate}


\section{Phase 2 Diverges After 200--500 Steps}

\begin{table}[H]
\centering
\caption{Phase~2 troubleshooting.}
\begin{tabularx}{\textwidth}{@{}lXX@{}}
\toprule
\textbf{Symptom} & \textbf{Cause} & \textbf{Fix} \\
\midrule
Floating-point exception &
  BCD scheme oscillating on strong gradients &
  Temporarily switch momentum to \textbf{Second Order Upwind} for
  1000 steps, then switch back to BCD \\
Max Courant $> 20$ &
  Time step too large for local cell size near junction &
  Reduce $\Delta t$ to \textbf{1e-4}; or reduce BOI element size
  from 5\,mm to 3\,mm (coarser mesh = larger stable $\Delta t$) \\
Species residual explodes &
  H$_2$ mass fraction going $> 1$ or $< 0$ &
  Enable \textbf{Bounded} on species under Solution Methods;
  reduce species URF to \textbf{0.7} \\
\bottomrule
\end{tabularx}
\end{table}


\section{Phase 3: DDES Causes Immediate Divergence}

The DDES switch activates LES-like behaviour in separated regions.
If the Phase~2 field has large RANS eddies that suddenly ``break up''
under DES, this can cause a spike:

\begin{enumerate}
  \item Reload \texttt{phase2\_transient\_rans.cas.h5}.
  \item Enable DDES.
  \item Reduce $\Delta t$ to \textbf{1e-4}\,s.
  \item Run 2000 steps at this reduced $\Delta t$.
  \item If stable, increase $\Delta t$ back to \textbf{2e-4}\,s.
\end{enumerate}


\section{General: ``Turbulent Viscosity Ratio Limited'' Warning}

This message appears when $\mu_t / \mu > 10^5$.  It usually means:

\begin{itemize}
  \item The turbulence model is generating unrealistically high
        $\mu_t$ at the junction stagnation point.
  \item \textbf{Not fatal} if it appears occasionally (especially in
        Phase~1).
  \item If it appears \textbf{every iteration}, reduce the turbulence
        URFs ($k$, $\omega$) to 0.5 and investigate the mesh in the
        junction region.
\end{itemize}


\section{General: Mass Imbalance $> 1\,\%$}

\begin{enumerate}
  \item Check that \textbf{all} boundary conditions are set correctly.
  \item Verify that the outlet backflow species fraction is \textbf{0}
        (not 1).
  \item Check for reversed flow at the outlet: \textbf{Flux Reports}
        $\rightarrow$ mass flow rate on outlet.  If negative, the
        outlet is too close or the pressure BC is wrong.
\end{enumerate}


% ════════════════════════════════════════════════════════════════
%  CHAPTER 9 — Quick-Reference Summary
% ════════════════════════════════════════════════════════════════
\chapter{Quick-Reference Summary}
\label{ch:summary}

\begin{table}[H]
\centering
\caption{Complete settings summary across all three phases.}
\small
\begin{tabularx}{\textwidth}{@{}lXXX@{}}
\toprule
\textbf{Setting} & \textbf{Phase 1} & \textbf{Phase 2} &
\textbf{Phase 3} \\
\midrule
Mode            & Steady   & Transient      & Transient \\
Turbulence      & SST RANS & SST RANS       & \textbf{SST-DDES} \\
P--V Coupling   & Coupled  & PISO           & PISO \\
Momentum Scheme & 2nd Upwind & \textbf{BCD} & BCD \\
Pressure Scheme & 2nd Order & PRESTO!       & PRESTO! \\
Transient Form. & ---      & Bounded 2nd Impl. & Bounded 2nd Impl. \\
$\Delta t$ (s)  & ---      & 5e-4 $\rightarrow$ 2e-4 & 2e-4 \\
Iterations      & 3000     & 5000 steps     & 25000 steps \\
Max Iter/Step   & ---      & 30 $\rightarrow$ 20 & 20 \\
\midrule
\multicolumn{4}{@{}l}{\textbf{Under-Relaxation / Controls}} \\
\midrule
Pressure        & CFL=50   & 0.3            & 0.3 \\
Momentum        & 0.5      & 0.7            & 0.7 \\
$k$ / $\omega$  & 0.7      & 0.8            & 0.8 \\
Species         & 0.8      & 0.9            & 0.9 \\
Energy          & 0.9      & 1.0            & 1.0 \\
\midrule
\multicolumn{4}{@{}l}{\textbf{Checkpoints}} \\
\midrule
Save as         & \texttt{phase1\_steady\_rans} &
\texttt{phase2\_transient\_rans} & \texttt{phase3\_ddes\_final} \\
\bottomrule
\end{tabularx}
\end{table}

\vfill

\noindent\fbox{\parbox{\dimexpr\textwidth-2\fboxsep-2\fboxrule}{%
\textbf{Pre-flight checklist before Phase~1:}
\begin{itemize}[nosep]
  \item[$\square$] Mesh quality checked (orthogonality $> 0.1$,
        skewness $< 0.95$)
  \item[$\square$] Named Selections present (main\_inlet,
        branch\_inlet, outlet, wall)
  \item[$\square$] Material = ch4-h2-mixture with ideal-gas density
  \item[$\square$] Operating pressure = 6\,900\,000\,Pa
  \item[$\square$] Operating density = 0
  \item[$\square$] Gravity enabled, $g_y = -9.81$
  \item[$\square$] Energy equation ON
  \item[$\square$] Species Transport ON, Reactions OFF
  \item[$\square$] $Sc_t = 0.5$ (set via TUI)
  \item[$\square$] Double precision enabled
\end{itemize}}}


\end{document}
